\section{Apex Gradients}
\label{sec:apex-gradients}

Apex gradients describe the maximal directional rate of change of the operator
spine energy functional. They identify the steepest ascent direction in the
80-dimensional embedding space and govern how NDH's internal geometry responds
to perturbation. Apex gradients interact with curvature-regulating operators,
including the 9D radial Laplacian, to maintain stability of the cognitive
manifold.

\subsection{Operator-Spine Energy Functional}

Let $\mathcal{E} : \mathbb{R}^{80} \to \mathbb{R}$ denote the operator-spine
energy functional defined over the 80-dimensional embedding space. The gradient
$\nabla \mathcal{E}(x)$ represents the local direction of maximal increase in
energy at state $x$. Apex gradients refine this notion by selecting the unit
direction that maximizes the directional derivative.

\subsection{Definition of Apex Gradient}

The apex gradient at state $x$ is defined by
\[
\nabla_{\mathrm{apex}} \mathcal{E}(x)
    = \operatorname*{arg\,max}_{\|v\| = 1}
      \left\langle \nabla \mathcal{E}(x), v \right\rangle.
\]
Since the inner product is maximized when $v$ aligns with the gradient, the
apex gradient reduces to the normalized gradient:
\[
\nabla_{\mathrm{apex}} \mathcal{E}(x)
    = \frac{\nabla \mathcal{E}(x)}{\|\nabla \mathcal{E}(x)\|}.
\]
This vector identifies the steepest ascent direction in the operator-spine
geometry.

\subsection{Role in Operator-Spine Dynamics}

Apex gradients determine the preferred direction of geometric deformation under
semantic excitation. They regulate:

\begin{itemize}
    \item stability of the 80D embedding manifold,
    \item reversibility of the $\Sigma/\Lambda$ save-state geometry,
    \item modal-boundary deformation behavior,
    \item spectral-recursion coupling alignment,
    \item holonomy-clamp tension.
\end{itemize}

Apex gradients therefore serve as the directional component of NDH's internal
physics, governing how the system responds to perturbation.

\subsection{Interaction with the 9D Radial Laplacian}

Let $f(r)$ be a radially symmetric component of the operator-spine energy
functional, where $r = \|x\|$ in $\mathbb{R}^9$. The 9D radial Laplacian is
given by
\[
\Delta_9 f(r)
    = f''(r) + \frac{8}{r} f'(r).
\]
This operator governs radial curvature diffusion in the operator spine.

Apex gradients determine the direction of maximal change, while the radial
Laplacian determines how curvature diffuses along radial directions. Their
interaction ensures that geometric deformation is both directionally coherent
and radially stable.

\subsection{Stability and Drift Neutrality}

Apex gradients contribute to drift-neutrality by aligning deformation with
stable directions in the operator-spine geometry. When combined with curvature
diffusion from the radial Laplacian, apex gradients prevent runaway deformation
and ensure that NDH maintains structural coherence across modal regions.

\subsection{Summary}

Apex gradients identify the steepest ascent direction of the operator-spine
energy functional and govern NDH's internal geometric response to perturbation.
Their interaction with the 9D radial Laplacian provides directional and radial
stability, ensuring coherence of the 80D embedding manifold and reversibility of
the $\Sigma/\Lambda$ save-state geometry.
